\documentclass[11pt,a4paper]{report}

% ==========================================
% 1. TYPOGRAPHIE & ENCODAGE
% ==========================================
\usepackage[utf8]{inputenc}
\usepackage[T1]{fontenc}
\usepackage[french]{babel}
\usepackage{mathpazo} % Police Palatino (corps)
\usepackage[scaled=0.95]{helvet} % Helvetica (titres)
\usepackage{microtype} 

% ==========================================
% 2. MISE EN PAGE & GEOMETRY
% ==========================================
\usepackage[a4paper, margin=2.5cm, top=3cm, bottom=3cm]{geometry}
\usepackage{parskip}
\usepackage{fancyhdr}
\usepackage{lastpage}

% ==========================================
% 3. COULEURS & DESIGN
% ==========================================
\usepackage[table,xcdraw]{xcolor}
\definecolor{mainblue}{RGB}{41, 57, 85}    % Bleu nuit profond
\definecolor{accentblue}{RGB}{70, 130, 180} % Bleu acier
\definecolor{lightgray}{RGB}{245, 245, 245} 
\definecolor{accentgreen}{RGB}{16,124,16}   % Vert succès

% ==========================================
% 4. TITRES (Simples & Robustes)
% ==========================================
\usepackage{titlesec}
\titleformat{\chapter}[display]
  {\normalfont\huge\bfseries\sffamily\color{mainblue}}
  {\chaptertitlename\ \thechapter}
  {20pt}
  {\Huge}
\titleformat{\section}
  {\Large\sffamily\bfseries\color{mainblue}}
  {\thesection}
  {1em}
  {}
  [{\titlerule[1pt]}]

% ==========================================
% 5. ÉLÉMENTS VISUELS (TCOLORBOX)
% ==========================================
\usepackage[most]{tcolorbox}

% Boîte Info
\newtcolorbox{infobox}[1][]{
  colback=lightgray, colframe=mainblue, fonttitle=\bfseries\sffamily,
  title=#1, boxrule=0.5pt, arc=2mm, left=4mm, right=4mm, top=3mm, bottom=3mm
}
% Boîte Résultat
\newtcolorbox{successbox}{
    colback=accentgreen!5, colframe=accentgreen, fonttitle=\bfseries,
    title=Résultat Clé, boxrule=1pt, rounded corners
}
% Boîte Équation
\newtcolorbox{eqbox}{
    colback=mainblue!5, colframe=mainblue, boxrule=1pt,
    left=5pt,right=5pt,top=5pt,bottom=5pt
}

% ==========================================
% 6. TABLEAUX & MATHS
% ==========================================
\usepackage{booktabs}
\usepackage{amsmath,amssymb}
\usepackage{colortbl}

% ==========================================
% 7. NAVIGATION
% ==========================================
\usepackage{hyperref}
\hypersetup{
    colorlinks=true,
    linkcolor=mainblue,
    urlcolor=accentblue
}

% CONFIGURATION DU HEADER
\pagestyle{fancy}
\fancyhf{}
\fancyhead[L]{\sffamily\small \textcolor{mainblue}{Maintenance Prédictive Deep RL}}
\fancyhead[R]{\sffamily\small \today}
\fancyfoot[C]{Page \thepage\ sur \pageref{LastPage}}
\renewcommand{\headrulewidth}{0.4pt}

% ==========================================
% DEBUT DU DOCUMENT
% ==========================================
\begin{document}

% --- PAGE DE GARDE PREMIUM ---
\begin{titlepage}
    \centering
    \vspace*{2cm}
    
    {\scshape\Large Rapport Technique Académique \par}
    \vspace{0.5cm}
    {\sffamily\bfseries\Huge\color{mainblue} Optimisation de la Maintenance Industrielle \par}
    \vspace{0.3cm}
    {\Large\itshape par Apprentissage par Renforcement Profond (Deep RL) \par}
    
    \vspace{2cm}
    
    \begin{tcolorbox}[colback=lightgray, colframe=white, arc=0mm, boxrule=0pt]
        \centering
        \vspace{0.5cm}
        Comparaison des stratégies \textbf{DQN} et \textbf{PPO} face aux méthodes traditionnelles.
        \vspace{0.5cm}
    \end{tcolorbox}

    \vspace{3cm}

    \begin{minipage}{0.4\textwidth}
        \begin{flushleft} \large
        \emph{Auteur :}\\
        Amine \textsc{AMLLAL}
        \end{flushleft}
    \end{minipage}
    \hfill
    \begin{minipage}{0.4\textwidth}
        \begin{flushright} \large
        \emph{Encadrant :}\\
        M. Tawfik \textsc{MASROUR}
        \end{flushright}
    \end{minipage}

    \vfill
    {\large Année Universitaire 2024-2025 \par}
    \vspace{1cm}
\end{titlepage}

\tableofcontents
\newpage

% --- CONTENU TECHNIQUE ENRICHI ---

\chapter{Introduction}
\section{Contexte et Enjeux}
La maintenance des équipements critiques représente un enjeu économique majeur. Les approches traditionnelles (réactive, préventive) peinent à trouver l'équilibre optimal entre le coût d'intervention et le risque de panne.
Ce projet propose une approche basée sur un agent intelligent autonome.

\begin{infobox}[Objectif Principal]
Développer un agent Deep RL capable de décider dynamiquement de la maintenance en maximisant le RUL (Remaining Useful Life) utilisé tout en minimisant les coûts.
\end{infobox}

\chapter{Formulation du Problème (MDP)}

\section{Définition Formelle}
Le problème est modélisé sous la forme d'un Processus de Décision Markovien $ \langle \mathcal{S}, \mathcal{A}, \mathcal{P}, \mathcal{R}, \gamma \rangle $.

\section{Espace d'États $\mathcal{S}$}
L'agent observe une fenêtre glissante de capteurs pour percevoir la dynamique temporelle de la dégradation.
\begin{eqbox}
\[ s_t = [x_{t-29}, x_{t-28}, \ldots, x_{t-1}, x_t] \in \mathbb{R}^{420} \]
\end{eqbox}
Où $x_t$ est un vecteur de 14 capteurs physiques (vibrations, température, pression...).

\section{Espace d'Actions $\mathcal{A}$}
À chaque pas de temps $t$, l'agent choisit une action binaire :
\begin{eqbox}
\[ \mathcal{A} = \{0 (\text{Continuer}), 1 (\text{Maintenance})\} \]
\end{eqbox}

\section{Fonction de Transition $\mathcal{P}$}
La dynamique du système est déterministe pour la dégradation :
\begin{eqbox}
\[
\text{RUL}_{t+1} = \begin{cases}
\text{RUL}_{\max} & \text{si } a_t = 1 \\
\max(0, \text{RUL}_t - \Delta_t) & \text{si } a_t = 0
\end{cases}
\]
\end{eqbox}
où la dégradation $\Delta_t$ est modélisée par :
\begin{eqbox}
\[
\Delta_t = \beta_{base} + \epsilon_t + \beta_{accel} \cdot \left(1 - \frac{\text{RUL}_t}{\text{RUL}_{\max}}\right)^{1.5}
\]
\end{eqbox}
avec $\beta_{base} = 1.0$, $\epsilon_t \sim \mathcal{N}(0, 2.5)$, et $\text{RUL}_{\max} = 125$ cycles.

\section{Fonction de Récompense $\mathcal{R}$ - DÉTAILLÉE}

\begin{infobox}[Composante Centrale]
La fonction équilibre trois objectifs critiques : maintenir la production, minimiser les coûts d'intervention et éviter absolument la panne.
\end{infobox}

\begin{eqbox}
\[
r_t = \begin{cases}
+1.0 & \text{si } a_t = 0 \text{ et RUL}_t > 0 \\
-(50 + 0.5 \cdot \text{RUL}_t) & \text{si } a_t = 1 \\
-200 & \text{si RUL}_t \leq 0
\end{cases}
\]
\end{eqbox}

Le coût de maintenance variable est défini par :
\[ C_{maint}(\text{RUL}_t) = C_{fixed} + \alpha \cdot \text{RUL}_t \]

\begin{table}[h]
\centering
\begin{tabular}{lcl}
\toprule
\rowcolor{mainblue} \color{white}\textbf{Paramètre} & \color{white}\textbf{Valeur} & \color{white}\textbf{Justification} \\ \midrule
$r_{running}$ & +1.0 & Bonus fonctionnement normal \\
$C_{fixed}$ & 50 € & Coût fixe de maintenance \\
$\alpha$ & 0.5 & Pénalité pour RUL gaspillé \\
$C_{failure}$ & 200 € & Coût de panne catastrophique \\ \bottomrule
\end{tabular}
\caption{Paramètres de la fonction de récompense}
\end{table}

\begin{successbox}
\textbf{Justification de la stratégie :}
\begin{enumerate}
    \item Le terme $\alpha \cdot \text{RUL}_t$ pénalise une maintenance trop précoce.
    \item $C_{failure} > C_{maint}$ encourage fortement la prévention.
    \item $r_{running}$ incite à maximiser la durée de vie opérationnelle.
\end{enumerate}
\end{successbox}

\chapter{Algorithmes et Implémentation}

\section{Deep Q-Network (DQN)}
DQN utilise un réseau de neurones pour approximer la fonction de valeur $Q(s, a)$.
\begin{itemize}
    \item \textbf{Architecture :} MLP (420 $\to$ 256 $\to$ 256 $\to$ 2).
    \item \textbf{Spécificité :} Utilisation d'un \emph{Replay Buffer} de 10,000 transitions.
\end{itemize}

\section{Proximal Policy Optimization (PPO)}
PPO est une méthode \emph{Policy-Gradient} qui optimise directement la politique $\pi_\theta$. Il s'est avéré plus stable pour ce problème de contrôle continu.
\begin{itemize}
    \item \textbf{Architecture :} Actor-Critic.
    \item \textbf{Avantage :} Mises à jour prudentes ("Clipping") évitant la divergence.
\end{itemize}

\chapter{Résultats et Comparaison}

\section{Performance des Algorithmes}

L'entraînement a été réalisé sur 100,000 pas de temps. PPO a montré une convergence plus rapide et plus stable que DQN.

\begin{table}[h]
\centering
\begin{tabular}{lcccc}
\toprule
\rowcolor{mainblue} \color{white}\textbf{Algorithme} & \color{white}\textbf{Coût Moyen} & \color{white}\textbf{Taux Panne} & \color{white}\textbf{RUL Utilisé} & \color{white}\textbf{Temps Entr.} \\ \midrule
DQN & 72 € & 3\% & 88\% & 120s \\
\textbf{PPO} & \textbf{65 €} & \textbf{2\%} & \textbf{92\%} & \textbf{95s} \\ \bottomrule
\end{tabular}
\caption{Comparaison PPO vs DQN (Moyenne sur 100 épisodes)}
\end{table}

\section{Comparaison avec les Méthodes Classiques}

L'agent PPO a ensuite été confronté aux stratégies industrielles standards sur une simulation longue durée (2000 cycles).

\begin{table}[h]
\centering
\begin{tabular}{lcccc}
\toprule
\rowcolor{mainblue} \color{white}\textbf{Stratégie} & \color{white}\textbf{Coût Total} & \color{white}\textbf{Panne} & \color{white}\textbf{RUL Utilisé} & \color{white}\textbf{Rang} \\ \midrule
Maintenance Périodique & 140 € & 5\% & 45\% & 3 \\
Maintenance à Seuils & 100 € & 2\% & 70\% & 2 \\
\rowcolor{lightgray} \textbf{IA (PPO)} & \textbf{65 €} & \textbf{2\%} & \textbf{92\%} & \textbf{1} \\ \bottomrule
\end{tabular}
\end{table}

\begin{successbox}
\textbf{Impact Industriel :}
L'approche par Deep RL (PPO) permet de \textbf{réduire les coûts de maintenance de 53\%} par rapport à une approche périodique, tout en maximisant l'utilisation de la durée de vie des composants (92\%).
\end{successbox}

\chapter{Conclusion et Perspectives}
Ce travail valide l'efficacité du Deep Reinforcement Learning pour la maintenance prédictive. L'agent autonome apprend une politique complexe qui surpasse les règles heuristiques humaines.
Les perspectives incluent le test sur des données réelles (NASA C-MAPSS) et l'intégration dans un système SCADA industriel.

\end{document}
